\documentclass[aspectratio=169,12pt]{beamer}
% Estilo canónico para las presentaciones de Métodos Numéricos 2026-2.
\usepackage[T1]{fontenc}
\usepackage[utf8]{inputenc}
\usepackage[spanish,es-nodecimaldot]{babel}
\usepackage{lmodern}
\usepackage{amsmath,amssymb,mathtools}
\usepackage{booktabs}
\usepackage{graphicx}
\usepackage{xcolor}
\usepackage{hyperref}

\definecolor{upazul}{HTML}{003B70}
\definecolor{upazulclaro}{HTML}{EAF2F8}
\definecolor{updorado}{HTML}{C5962E}
\definecolor{upgris}{HTML}{4A5560}
\definecolor{upverde}{HTML}{167C5A}
\definecolor{uprojo}{HTML}{B33A3A}

\setbeamercolor{normal text}{fg=black,bg=white}
\setbeamercolor{structure}{fg=upazul}
\setbeamercolor{frametitle}{fg=white,bg=upazul}
\setbeamercolor{title}{fg=upazul}
\setbeamercolor{subtitle}{fg=upgris}
\setbeamercolor{block title}{fg=white,bg=upazul}
\setbeamercolor{block body}{fg=black,bg=upazulclaro}
\setbeamercolor{alerted text}{fg=uprojo}
\setbeamercolor{example text}{fg=upverde}

\setbeamerfont{title}{size=\Huge,series=\bfseries}
\setbeamerfont{subtitle}{size=\Large}
\setbeamerfont{frametitle}{size=\LARGE,series=\bfseries}
\setbeamerfont{normal text}{size=\large}
\setbeamerfont{block title}{series=\bfseries}

\setbeamertemplate{navigation symbols}{}
\setbeamertemplate{itemize item}{\color{updorado}$\blacktriangleright$}
\setbeamertemplate{itemize subitem}{\color{upazul}$\bullet$}
\setbeamertemplate{blocks}[rounded][shadow=false]

\setbeamertemplate{frametitle}{%
  \nointerlineskip
  \begin{beamercolorbox}[wd=\paperwidth,ht=0.78cm,dp=0.24cm,leftskip=0.55cm]{frametitle}
    \usebeamerfont{frametitle}\insertframetitle
  \end{beamercolorbox}}

\setbeamertemplate{footline}{%
  \leavevmode
  \hbox{%
    \begin{beamercolorbox}[wd=.82\paperwidth,ht=2.7ex,dp=1.2ex,leftskip=.5cm]{author in head/foot}
      \color{upgris}\scriptsize Métodos Numéricos 2026-2 \hspace{0.8em} | \hspace{0.8em} Universidad del Pacífico
    \end{beamercolorbox}%
    \begin{beamercolorbox}[wd=.18\paperwidth,ht=2.7ex,dp=1.2ex,rightskip=.5cm plus1fil]{date in head/foot}
      \color{upazul}\scriptsize\hfill\insertframenumber/\inserttotalframenumber
    \end{beamercolorbox}}
  \vskip0pt}

\newcommand{\concepto}[1]{\textcolor{upazul}{\textbf{#1}}}
\newcommand{\resultado}[1]{\textcolor{upverde}{\textbf{#1}}}
\newcommand{\advertencia}[1]{\textcolor{uprojo}{\textbf{#1}}}
\newcommand{\codigo}[1]{\colorbox{upazulclaro}{\strut\texttt{#1}}}
\newcommand{\fuente}[1]{\vfill\hfill{\tiny\color{upgris}#1}}


\usepackage{listings}
\lstset{language=Python,basicstyle=\ttfamily\footnotesize,columns=fullflexible,
keepspaces=true,showstringspaces=false,breaklines=true,
literate={á}{{\'a}}1 {é}{{\'e}}1 {í}{{\'i}}1 {ó}{{\'o}}1 {ú}{{\'u}}1 {ñ}{{\~n}}1}
\newcommand{\R}{\mathbb{R}}
\newcommand{\norm}[1]{\left\lVert#1\right\rVert}
\newcommand{\PiN}{\mathcal P}
\newtheorem{thm}{Teorema}
\newtheorem{lema}{Lema}
\newtheorem{defi}{Definición}
\author{Oswaldo José Velásquez Castañón}
\institute{Universidad del Pacífico}
\date{Métodos Numéricos -- 2026-2}
\title{Integración numérica}
\subtitle{Cuadratura de Gauss y esperanzas}
\begin{document}
\begin{frame}[plain]
\begin{columns}[T,onlytextwidth]
\begin{column}{.70\textwidth}
\vspace{.7cm}{\usebeamerfont{title}\color{upazul}\inserttitle\par}
\vspace{.35cm}{\usebeamerfont{subtitle}\color{upgris}\insertsubtitle\par}
\vspace{1cm}{\large\insertauthor\par}{\normalsize\insertdate\par}
\end{column}
\begin{column}{.25\textwidth}\vspace{.55cm}\includegraphics[width=\linewidth]{UPlogo.pdf}\end{column}
\end{columns}
\end{frame}
\begin{frame}{Ruta de la clase}
\begin{itemize}
\item Elegir nodos para alcanzar la mayor exactitud polinomial posible.
\item Construir Gauss--Legendre y justificar sus pesos positivos.
\item Calcular nodos y pesos con Golub--Welsch.
\item Usar el peso adecuado para una esperanza normal o exponencial.
\end{itemize}
El complemento final conserva las reglas de integración doble del material anterior.
\end{frame}

\begin{frame}{La elección de los nodos}
Con $m$ nodos prefijados, la integración del interpolante asegura exactitud
hasta grado $m-1$. Ahora elegimos también los nodos:
\[I_w(f)=\int_a^b w(x)f(x)\,dx\approx
G_m(f)=\sum_{j=1}^m\lambda_j f(x_j).\]
Suponemos un peso integrable $w\geq0$, positivo casi en todas partes,
y $a<b$ finitos para los teoremas que siguen.
\begin{defi}
La regla es de tipo Gauss si es exacta para todo polinomio de grado
$\leq2m-1$.
\end{defi}
\end{frame}

\begin{frame}{Grado máximo de exactitud}
\begin{thm}
Ninguna regla con $m$ nodos distintos puede ser exacta para todos los
polinomios de grado $\leq2m$ bajo las hipótesis anteriores.
\end{thm}
\textbf{Demostración.} Tome
\[q(x)=\prod_{j=1}^m(x-x_j)^2.\]
Entonces $G_m(q)=0$, mientras que $I_w(q)>0$. Como $\deg q=2m$,
la exactitud en todo $\PiN_{2m}$ es imposible.
Una regla puede integrar algunos polinomios de grado mayor por simetría;
la afirmación se refiere a todos los polinomios hasta un grado dado.
\end{frame}

\begin{frame}{Polinomios ortogonales}
Definimos
\[\langle p,q\rangle_w=\int_a^b w(x)p(x)q(x)\,dx.\]
Al ortonormalizar $1,x,x^2,\ldots$ obtenemos $p_0,p_1,\ldots$, con
\[\deg p_k=k,\quad \langle p_i,p_j\rangle_w=\delta_{ij}.\]
\begin{thm}
Los $m$ ceros de $p_m$ son reales, simples y pertenecen a $(a,b)$.
\end{thm}
Esos ceros serán los nodos de la regla de Gauss.
\end{frame}

\begin{frame}{Ceros del polinomio ortogonal}
Sean $\xi_1,\ldots,\xi_k$ los ceros donde $p_m$ cambia de signo en $(a,b)$,
y tome $q(x)=\prod_{j=1}^k(x-\xi_j)$.
El producto $p_mq$ conserva signo y no es idénticamente nulo; por ello
\[\langle p_m,q\rangle_w\ne0.\]
Si $k<m$, entonces $q\in\PiN_{m-1}$ y la ortogonalidad daría
$\langle p_m,q\rangle_w=0$, contradicción.
Luego $k=m$. Un polinomio de grado $m$ con $m$ cambios de signo interiores
tiene exactamente $m$ raíces simples, todas en $(a,b)$.
\end{frame}

\begin{frame}{Teorema de construcción de Gauss}
\begin{thm}
Sean $x_1,\ldots,x_m$ los ceros de $p_m$, y $L_j$ las bases de Lagrange
para esos nodos. Los pesos
\[\lambda_j=\int_a^b w(x)L_j(x)\,dx\]
definen una regla exacta en $\PiN_{2m-1}$.
\end{thm}
\textbf{Demostración.} Si $p\in\PiN_{2m-1}$, divida
$p=qp_m+r$, con $\deg q,\deg r\leq m-1$.
Por ortogonalidad $I_w(qp_m)=0$. En los nodos $p(x_j)=r(x_j)$.
La regla interpolatoria integra $r$, de modo que $G_m(p)=I_w(r)=I_w(p)$.
\end{frame}

\begin{frame}{Positividad y suma de los pesos}
\begin{thm}
Los pesos de Gauss son positivos y satisfacen
\[\lambda_j>0,\qquad\sum_{j=1}^m\lambda_j=\int_a^b w(x)\,dx=\mu_0.\]
\end{thm}
\textbf{Demostración.} Como $L_j^2\in\PiN_{2m-2}$, la exactitud implica
\[\lambda_j=G_m(L_j^2)=I_w(L_j^2)>0.\]
La suma se obtiene aplicando la regla a la constante $1$.
Si $|\delta_j|\leq\eta$, entonces $|\Delta G_m|\leq\mu_0\eta$.
\end{frame}

\begin{frame}{Teorema del error de Gauss}
\begin{thm}
Sea $\omega_m(x)=\prod_{j=1}^m(x-x_j)$ el polinomio mónico de los nodos.
Si $f\in C^{2m}[a,b]$, existe $\xi\in(a,b)$ tal que
\[I_w(f)-G_m(f)=\frac{f^{(2m)}(\xi)}{(2m)!}
\int_a^b w(x)\omega_m(x)^2\,dx.\]
\end{thm}
El factor integral depende del peso, del intervalo y de los nodos.
No se omite al construir una cota.
La fórmula se obtiene integrando el error del interpolante de Hermite que
iguala valores y derivadas en los nodos, de grado $\leq2m-1$.
\end{frame}

\begin{frame}{Gauss--Legendre con dos nodos}
En $[-1,1]$, el peso es $w=1$. Por simetría buscamos nodos $\pm z$
y pesos iguales $\lambda$.
\[2\lambda=\int_{-1}^1 1\,dx=2,\qquad
2\lambda z^2=\int_{-1}^1x^2\,dx=\frac23.\]
Por tanto,
\[\boxed{G_2(f)=f(-1/\sqrt3)+f(1/\sqrt3).}\]
Los momentos impares se anulan por simetría. La regla es exacta hasta grado $3$.
Por ejemplo, $G_2(x^4)=2/9\ne2/5$.
\end{frame}

\begin{frame}{Transformación a un intervalo general}
Con
\[x=\frac{a+b}{2}+\frac{b-a}{2}z,\qquad dx=\frac{b-a}{2}\,dz,\]
se obtiene
\[\int_a^b f(x)\,dx\approx\frac{b-a}{2}
\sum_{j=1}^m\lambda_j f\left(\frac{a+b}{2}+\frac{b-a}{2}z_j\right).\]
Los $z_j,\lambda_j$ se calculan para $[-1,1]$.
En el intervalo original, los pesos suman $b-a$.
Olvidar el factor $(b-a)/2$ cambia la integral que se está calculando.
\end{frame}

\begin{frame}{El mismo excedente, ahora con Gauss}
Para $p(q)=50/(1+q/10)$, $I=500\log3$ y $EC=I-1000/3$:
\begin{center}\begin{tabular}{rrrr}\toprule
Nodos $m$&$G_m$&$|I-G_m|$&$EC$ aproximado\\\midrule
2&545.4545454545&3.8515988795&212.1212121212\\
4&549.2851768247&0.0209675094&215.9518434913\\
8&549.3061437596&$5.745\times10^{-7}$&215.9728104263\\
12&549.3061443340&$\approx10^{-11}$&215.9728110007\\\bottomrule
\end{tabular}\end{center}
El costo básico es $m$ evaluaciones de $p$.
La comparación con Newton--Cotes debe contar evaluaciones, no solo el parámetro de la regla.
\end{frame}

\begin{frame}[fragile]{Gauss--Legendre en NumPy}
\begin{lstlisting}
import numpy as np
from numpy.polynomial.legendre import leggauss

def gauss_legendre(f, a, b, m):
    z, w = leggauss(m)
    x = (a+b)/2 + (b-a)/2*z
    return (b-a)/2*np.dot(w, f(x))

p = lambda q: 50/(1+q/10)
I = gauss_legendre(p, 0., 20., 8)
EC = I - 20*p(20.)
\end{lstlisting}
Los nodos y pesos se pueden reutilizar para otras funciones en el mismo intervalo.
\end{frame}

\begin{frame}[fragile]{Verificación por momentos}
\begin{lstlisting}
z, w = leggauss(4)
assert np.all(w > 0)
assert np.isclose(w.sum(), 2.)
for k in range(8):
    exacta = 0. if k % 2 else 2/(k+1)
    assert abs(np.dot(w, z**k)-exacta) < 1e-12
\end{lstlisting}

Estas pruebas verifican constantes, momentos pares y simetría.
Una prueba con una sola función no comprueba el grado completo de exactitud.
\end{frame}

\begin{frame}{Golub--Welsch: matriz de Jacobi}
Para polinomios ortonormales con coeficientes principales positivos,
\[x p_k(x)=b_{k+1}p_{k+1}(x)+a_kp_k(x)+b_kp_{k-1}(x),\qquad b_k>0.\]
Se forma la matriz simétrica tridiagonal
\[J_m=\begin{pmatrix}
a_0&b_1&&\\b_1&a_1&\ddots&\\&\ddots&\ddots&b_{m-1}\\&&b_{m-1}&a_{m-1}
\end{pmatrix}.\]
\begin{block}{Nodos y pesos}
Si $J_mv_j=z_jv_j$ y $\norm{v_j}_2=1$, entonces
$x_j=z_j$ y $\lambda_j=\mu_0(v_j)_1^2$.
\end{block}
\end{frame}

\begin{frame}{Golub--Welsch para Legendre}
En este caso,
\[a_k=0,\qquad b_k=\frac{k}{\sqrt{(2k-1)(2k+1)}},\qquad\mu_0=2.\]
La recurrencia relaciona el determinante de $xI-J_m$ con el polinomio ortogonal
mónico: sus ceros son los valores propios.
Los pesos se recuperan de la primera componente de cada vector propio.
\begin{block}{Conexión con sistemas lineales}
La cuadratura conduce a un problema de valores y vectores propios de una
matriz tridiagonal simétrica. Se aprovecha esa estructura.
\end{block}
\end{frame}

\begin{frame}[fragile]{Implementación de Golub--Welsch}
\begin{lstlisting}
from scipy.linalg import eigh_tridiagonal

def golub_welsch_legendre(m):
    k = np.arange(1., m)
    d = np.zeros(m)
    e = k/np.sqrt(4*k*k-1)
    z, V = eigh_tridiagonal(d, e)
    w = 2*V[0, :]**2
    return z, w

z, w = golub_welsch_legendre(8)
\end{lstlisting}
Compare con \texttt{leggauss(8)}, verificando nodos y pesos, no el signo de los vectores propios.
\end{frame}

\begin{frame}{Cuadratura para una esperanza normal}
Si $Z\sim N(0,1)$, una regla de Gauss--Hermite puede escribirse como
\[\mathbb E[g(Z)]\approx\sum_{j=1}^m p_jg(z_j),\qquad\sum_jp_j=1.\]
Para $X=\mu+\sigma Z$,
\[\mathbb E[g(X)]\approx\sum_jp_jg(\mu+\sigma z_j).\]
En esta normalización, la matriz de Jacobi tiene diagonal cero y
fuera de diagonal $b_k=\sqrt{k}$.
Los pesos son $p_j=(v_j)_1^2$ porque la masa del peso normalizado es $1$.
\end{frame}

\begin{frame}{Polinomios de Hermite normalizados}
Bajo el producto interno $\langle f,g\rangle=\mathbb E[f(Z)g(Z)]$,
\[H_{-1}=0,\quad H_0=1,\quad
H_{k+1}(z)=\frac z{\sqrt{k+1}}H_k(z)-\sqrt{\frac{k}{k+1}}H_{k-1}(z).\]
Se cumple $H_k'(z)=\sqrt{k}H_{k-1}(z)$.
Si los $z_j$ son los ceros de $H_m$, los pesos normalizados son
\[p_j=\frac1{mH_{m-1}(z_j)^2},\qquad\sum_jp_j=1.\]
Estas fórmulas corresponden a la normal estándar, no al peso $e^{-z^2}$.
\end{frame}

\begin{frame}{Dos convenciones de Gauss--Hermite}
\small
\begin{center}\begin{tabular}{lll}\toprule
Rutina de NumPy&Peso integrado&Conversión a normal estándar\\\midrule
\texttt{hermgauss}&$e^{-z^2}$&$Z=\sqrt2z$, pesos $w/\sqrt\pi$\\
\texttt{hermegauss}&$e^{-z^2/2}$&$Z=z$, pesos $w/\sqrt{2\pi}$\\\bottomrule
\end{tabular}\end{center}
\[\mathbb E[g(Z)]\approx\frac1{\sqrt{2\pi}}\sum_jw_jg(z_j)
\quad\text{con \texttt{hermegauss}}.\]
Compruebe $\sum p_j=1$, $\sum p_jz_j=0$ y $\sum p_jz_j^2=1$.
Mezclar las convenciones modifica la varianza de la distribución integrada.
\end{frame}

\begin{frame}{Utilidad con consumo lognormal}
Suponga $C=\exp(\mu+\sigma Z)$, $Z\sim N(0,1)$, $\mu=0.5$, $\sigma=0.4$.
Para $u(C)=-1/C$,
\[\mathbb E[u(C)]=-\exp(-\mu+\sigma^2/2)=-0.657046819815.\]
\begin{center}\begin{tabular}{rrr}\toprule
Nodos&Gauss--Hermite&Error absoluto\\\midrule
2&$-0.655703538888$&$1.3433\times10^{-3}$\\
4&$-0.657046574205$&$2.4561\times10^{-7}$\\
8&$-0.657046819815$&$\approx3\times10^{-16}$\\\bottomrule
\end{tabular}\end{center}
Es un modelo didáctico. La expresión exacta permite auditar la integración.
\end{frame}

\begin{frame}[fragile]{Esperanza normal en Python}
\begin{lstlisting}
from numpy.polynomial.hermite_e import hermegauss
mu, sigma = .5, .4
z, w = hermegauss(8)
pesos = w/np.sqrt(2*np.pi)
utilidad = -np.exp(-(mu + sigma*z))
EU = np.dot(pesos, utilidad)
assert np.isclose(pesos.sum(), 1.)
assert np.isclose(np.dot(pesos, z*z), 1.)
\end{lstlisting}
La regla de ocho nodos se reutilizará al compararla con Montecarlo en la siguiente clase.
\end{frame}

\begin{frame}{Complemento: Gauss--Chebyshev}
Para el peso $w(z)=1/\sqrt{1-z^2}$ en $(-1,1)$,
\[\int_{-1}^1\frac{g(z)}{\sqrt{1-z^2}}\,dz
\approx\frac\pi m\sum_{j=1}^m g(z_j),\qquad
z_j=\cos\frac{(2j-1)\pi}{2m}.\]
Se conocen explícitamente los nodos y pesos.
Para una integral sin ese peso, tras el cambio afín $x=\phi(z)$ se usaría
\[g(z)=f(\phi(z))\sqrt{1-z^2}.\]
Ese factor puede reducir la regularidad en los extremos; conviene elegir
una familia cuyo peso coincida con el problema.
\end{frame}

\begin{frame}{Complemento: Gauss--Laguerre}
Para el peso $e^{-z}$ en $[0,\infty)$,
\[\int_0^\infty e^{-z}g(z)\,dz\approx\sum_jw_jg(z_j).\]
La matriz de Jacobi tiene diagonal $1,3,5,\ldots,2m-1$
y fuera de diagonal $1,2,\ldots,m-1$.
Si $X\sim\operatorname{Exp}(\lambda)$, el cambio $z=\lambda x$ da
\[\mathbb E[g(X)]=\int_0^\infty e^{-z}g(z/\lambda)\,dz
\approx\sum_jw_jg(z_j/\lambda).\]
El cambio de escala evita introducir factores exponenciales adicionales.
\end{frame}

\begin{frame}{Complemento: otras formas de Laguerre}
La misma cuadratura puede usarse en las identidades
\[\lambda\int_a^\infty e^{-\lambda x}f(x)\,dx
\approx\lambda e^{-\lambda a}\sum_jw_je^{(1-\lambda)z_j}f(z_j+a),\]
\[\int_0^\infty f(x)\,dx\approx\sum_jw_je^{z_j}f(z_j).\]
Se requieren integrales convergentes y valores numéricos manejables.
Las reponderaciones con exponenciales grandes pueden ser inestables;
si hay una distribución conocida, se prefiere transformar a su peso natural.
\end{frame}

\begin{frame}{Integración en dos dimensiones}
La cubicación aproxima
\[\iint_R f(x,y)\,dA.\]
En una región rectangular, se puede integrar el interpolante bilineal de
cada celda. En una región triangulada, se integra el interpolante afín.
Estas construcciones son las versiones en dos dimensiones de integrar
el interpolante de una variable.
También se pueden formar productos de reglas de Gauss de una dimensión.
\end{frame}

\begin{frame}{Trapecio en una celda rectangular}
Para $R_{ij}=[x_i,x_{i+1}]\times[y_j,y_{j+1}]$, de lados $h,k$,
\[\iint_{R_{ij}}f(x,y)\,dA\approx\frac{hk}{4}
[f_{ij}+f_{i+1,j}+f_{i+1,j+1}+f_{i,j+1}].\]
Es la integral exacta del interpolante bilineal de sus cuatro vértices.
Al sumar celdas en una malla rectangular uniforme,
\[Q(f)=hk\sum_{i=0}^n\sum_{j=0}^m a_i b_j f_{ij},\]
donde $a_0=a_n=b_0=b_m=1/2$ y los demás coeficientes valen $1$.
\end{frame}

\begin{frame}{Pesos de la regla rectangular compuesta}
La forma equivalente del material anterior es
\[Q(f)=\frac{hk}{4}\left[
\sum_{\text{esquinas}}f_{ij}+2\sum_{\text{bordes sin esquinas}}f_{ij}
+4\sum_{\text{interiores}}f_{ij}\right].\]
Las esquinas pertenecen a una celda, los puntos de borde a dos
y los interiores a cuatro.
La suma de pesos debe ser $(b-a)(d-c)$.
Para una regla de Gauss producto con $m$ nodos por eje, se necesitan $m^2$
evaluaciones; en dimensión $d$, $m^d$.
\end{frame}

\begin{frame}{Región triangulada}
Sea un triángulo con vértices $p_i,p_j,p_k$, área $A_\Delta$ y valores $f_i,f_j,f_k$.
La integral del interpolante afín es
\[Q_\Delta(f)=\frac{A_\Delta}{3}(f_i+f_j+f_k).\]
Si
\[D=\det\begin{pmatrix}1&x_i&y_i\\1&x_j&y_j\\1&x_k&y_k\end{pmatrix},\]
entonces $A_\Delta=|D|/2$ y $Q_\Delta=|D|(f_i+f_j+f_k)/6$.
La fórmula es exacta para funciones afines; la integral de una región poligonal
se aproxima sumando los triángulos de la malla.
\end{frame}

\begin{frame}{Actividad y fuentes}
\small
\begin{enumerate}
\item Compare ocho nodos de Gauss con trapecio y Simpson a costo similar.
\item Compruebe la suma de pesos antes y después de transformar el intervalo.
\item Calcule una esperanza normal con ambas convenciones de Hermite.
\end{enumerate}
\vfill
Material del curso: capítulo 7 y notas \textit{Cuadratura Gaussiana en la Práctica}
(atribuidas a Paul Klein en \texttt{golub.tex}).
\href{https://dlmf.nist.gov/3.5}{NIST DLMF 3.5}; Golub y Welsch (1969).
Documentación: \href{https://numpy.org/doc/stable/reference/generated/numpy.polynomial.legendre.leggauss.html}{leggauss},
\href{https://numpy.org/doc/stable/reference/generated/numpy.polynomial.hermite_e.hermegauss.html}{hermegauss},
\href{https://docs.scipy.org/doc/scipy/reference/generated/scipy.linalg.eigh_tridiagonal.html}{eigh\_tridiagonal}.
\end{frame}

\end{document}
